\section{Entropy profiles and their tangents}\label{sec:entropy}

We retain the calibrated path constructed in Section~\ref{sec:geo}.
At its differentiability point $y_0$, write
\[
x_0=x(y_0),\qquad d_0=d(y_0),\qquad
a=-x'(y_0),\qquad b=-(x+d)'(y_0),\qquad \lambda=a-b.
\]
Thus $a\geq0$, $b\leq1$, and the path widths, on the scale
$y=y_0+\delta v$, have depths
$x_0-\delta av+o(\delta)$ and
$x_0+d_0-\delta bv+o(\delta)$.
The first-order logarithmic unit is
\[
L_{\log}=\log\mathcal M=\delta\log M\longrightarrow\infty.
\]
All the errors from the path construction are $o(L_{\log})$ in
logarithmic units.  The stretched interval in
Lemma~\ref{geo:calibration} contains every fixed bounded interval of
offsets in the limit.  In particular, choosing an earlier or later
\emph{fixed} offset will not leave the available path.

The task in this section is to preserve the cap, heavy-list, and
axis-prediction information in finite conditional laws.  The profile
rules used later come from those laws and their hereditary
regularization.

We use discrete Shannon entropy, denoted by $\Ent$, with natural
logarithms; see Shannon~\cite{Shannon1948}.  The standard chain rule,
conditioning, and submodularity rules are collected, for example,
in \cite[arXiv:0906.4387v5, Appendix~A]{Tao2010Entropy}.
An expression such as $\Ent(Z\mid\mathcal B)$ is a conditional entropy; the base label $\mathcal B$ need not be discrete.
It determines $P=(U,V)$.  Every measured variable below is discretized,
and its range at a finite stage has at most $\exp(O(\log M))$ elements.

\subsection{Admissible sheared states}

Suppose first that $\lambda\ne0$, so $d_0>0$.  Let $R^\circ$ be the
reference slope chosen for each particle and used throughout the
calibrated local epochs in Lemma~\ref{geo:calibration}, in a fixed
bounded slope chart.  For offsets $z=(C,B,D,A,Y,R)$, put
$t=T_{y_0+\delta Y}$ and $r=[R^\circ]_{d_0+\delta R}$.
The state $Z(z)$ is the ordered tuple of bins of the following
measured coordinates, at the corresponding absolute depths:
\begin{equation}\label{eq:source-9}
\begin{array}{c|c|c}
\text{offset}&\text{measured coordinate}&\text{absolute depth}\\ \hline
C&U_1&x_0+\delta C\\
B&U_2-rU_1&x_0+d_0+\delta B\\
D&(X_t)_1&x_0+y_0+\delta D\\
A&(X_t)_2-r(X_t)_1&x_0+y_0+d_0+\delta A\\
Y&t&y_0+\delta Y\\
R&r&d_0+\delta R
\end{array}
\end{equation}
The notation for a depth includes the
dyadic rounding convention of Section~\ref{sec:geo}.
The domain of admissible offsets is the closed convex set
\begin{equation}\label{eq:source-10}
\Omega=\{B\leq C+R,\quad D\leq C+Y,\quad
                  A\leq B+Y,\quad A\leq D+R\}.
\end{equation}
Full alignment means equality in all four constraints.  It can be
parametrized by $(C,Y,R)$, with
$B=C+R$, $D=C+Y$, and $A=C+Y+R$.

We write $z\leq z'$ when every offset of $z$ is at most the
corresponding offset of $z'$.  The meet $z\wedge z'$ and join
$z\vee z'$ are their coordinatewise minimum and maximum.
A comparison rectangle has corners $z\wedge z'$, $z$, $z'$,
and $z\vee z'$; all four must belong to $\Omega$.

\begin{lemma}[Finite-state comparison]\label{ent:finite-states}
On any fixed bounded subset of $\Omega$, comparable states are nested
up to bounded ambiguity.  A refinement costs at most the logarithm of
the product of its scalar resolution ratios, up to an additive
constant.  For a coordinate rectangle whose four corners belong to
$\Omega$, the two side states together determine the upper state up
to bounded ambiguity, and each determines the bottom state up to
bounded ambiguity.  Consequently their entropies satisfy monotonicity,
the scalar-coordinate Lipschitz bounds, and submodularity, with
additive bounded errors.

If only the parameter cutoffs increase, the entropy increment is at
least the corresponding parameter entropy increment conditional on
$\mathcal B$, again up to a bounded error.
\end{lemma}

\begin{proof}
Write the spatial entries in the order of
\eqref{eq:source-9} as $(x,h,z,w)$.  Changing the truncated parameters
by $(\Delta t,\Delta r)$ gives the exact update
\begin{equation}\label{ent:shear-update}
(x,h,z,w)\longmapsto
\bigl(x,h-\Delta r\,x,z+\Delta t\,x,
w+\Delta t\,h-\Delta r\,z-\Delta t\Delta r\,x\bigr).
\end{equation}
For component widths $\beta_x,\beta_h,\beta_z,\beta_w$ and parameter
widths $\beta_t,\beta_r$, the errors in this update at the starting
component resolutions are bounded precisely when
\[
\beta_x\beta_r\lesssim\beta_h,\qquad
\beta_x\beta_t\lesssim\beta_z,\qquad
\beta_t\beta_h\lesssim\beta_w,\qquad
\beta_r\beta_z\lesssim\beta_w.
\]
These are \eqref{eq:source-10}, including the absolute background
depths.  They also bound the doubly sheared error, since
$\beta_t\beta_r\beta_x\lesssim\beta_t\beta_h\lesssim\beta_w$.
The shifts of bin centers are known from the parameter labels and
can be subtracted exactly.  Only the within-bin errors require these
inequalities; large absolute center coordinates cause no additional
loss.

A finer state supplies finer parameter bins and finer spatial
measurements.  Applying \eqref{ent:shear-update} recovers a coarser
state with a bounded number of possible bin labels.  Conversely,
inside a coarse state, each newly tested scalar coordinate and each
parameter has at most a constant times its resolution ratio many
refinements.  This proves the first assertions.

For a coordinate rectangle, the side states supply the finer pair of
parameter labels.  Each finest spatial component is already measured
by at least one side.  Update that side to the finer parameter pair at
its own component accuracy, using the inequalities valid at that
side.  This recovers the upper state.  The same update proves that
either side recovers the bottom state.  Applying the usual entropy
submodularity inequality to the two side labels proves the asserted
inequality with an additive bounded error.

Finally let $Q_0,Q_1$ be the coarse and fine parameter pairs in a
parameter-only refinement, and let $Z_0,Z_1$ be its states.  The
nesting just proved gives
\begin{align*}
\Ent(Z_1)-\Ent(Z_0)
&\geq \Ent(Q_1\mid Z_0)-O(1)\\
&\geq \Ent(Q_1\mid Z_0,\mathcal B)-O(1)\\
&=\Ent(Q_1\mid Q_0,\mathcal B)-O(1).
\end{align*}
For the last equality, $Z_0$ is determined by $\mathcal B$ and $Q_0$,
whereas $Z_0$ includes $Q_0$.  No independence between the two
parameters, or between either parameter and $\mathcal B$, is used.
\end{proof}

\subsection{Regularization before taking a profile}

Entropy values alone do not give the probability bounds needed in a
projection test.  We first control finite support sizes and maximal
cell weights uniformly under the restrictions that follow; only then
do we pass to an entropy profile.  A subpower restriction here means
a restriction of relative mass $\exp(-o(L_{\log}))$.

\begin{lemma}[Hereditary regularization]\label{ent:regularization}
For any fixed finite family of states, their joint labels, and
conditional parameter labels given $\mathcal B$, one may restrict to
a subpower-density law with the following properties.  At every
unconditional state of absolute entropy $h$,
\begin{equation}\label{ent:absolute-flatness}
\#\supp Z\leq\exp\bigl(h+o(L_{\log})\bigr),\qquad
\max_z\Pr(Z=z)\leq\exp\bigl(-h+o(L_{\log})\bigr).
\end{equation}
Analogous conditional bounds hold on the retained base-label fibers.
For comparable states of absolute entropies $h_0,h_1$, typical coarse
bins have fine support at most
$\exp(h_1-h_0+o(L_{\log}))$ and maximal conditional fine weight at
most $\exp(-(h_1-h_0)+o(L_{\log}))$.

These assertions can hold simultaneously on any prescribed
countable family of successively finer finite grids, interpreted
diagonally.  Subsequent subpower restrictions preserve the referenced
entropy profiles and these estimates, after discarding exceptional
parent or base-label fibers.  If lower weights are required in a
finite collection of tests, bins of smaller weight than the
corresponding reciprocal-support exponent can be discarded with
subpower slack; their retained lower weights refer, in particular, to
the law before that last discard.
\end{lemma}

\begin{proof}
Fix a finite family first.  The number of bins in each of its labels
is at most $\exp(O(\log M))$.  Tails of probabilities smaller than an
appropriately large negative power of $M$ have negligible total mass,
by this count.  Partition the remaining probabilities into dyadic
classes.  There are $O(\log M)$ relevant classes per label.  Do the
same for all the prescribed joint labels and, pointwise on the base
label, for the prescribed conditional parameter probabilities.
Pigeonholing selects common classes on an event of relative mass
\[
\exp\bigl(-O_J(\log\log M)\bigr)
       =\exp\bigl(-o(L_{\log})\bigr),
\]
where $J$ is the fixed number of tests.  We choose all diagonal
parameters slowly enough that this remains true as the grids grow.

Suppose a selected unconditional class has original weight comparable
to $w$.  It contains at most $O(w^{-1})$ bins.  After restriction and
normalization by a subpower density, its maximal weight is at most
$w\exp(o(L_{\log}))$.  Its new absolute entropy is therefore
$-\log w+o(L_{\log})$, by the upper support bound and the lower bound
on entropy supplied by the maximal atom.  This proves
\eqref{ent:absolute-flatness}.  For conditional classes, discard
base-label fibers on which the relative retained mass is smaller
than a sufficiently slowly vanishing power.  Such fibers have
negligible total retained mass.  The identical support and maximal
weight argument then applies on every remaining fiber.  This
procedure is also valid when $\mathcal B$ takes values in a standard
Borel space: the dyadic class of a conditional probability is a
measurable finite label, and no count of base labels is involved.

Here is the conditional-parent deduction explicitly.  Replace the
coarse and fine labels by their joint label when necessary.  By
Lemma~\ref{ent:finite-states}, this changes the fine support and
entropy bounds by bounded factors and additive constants.  Write the
errors in \eqref{ent:absolute-flatness} as $\varepsilon L_{\log}$.
If a parent has more than
$\exp(h_1-h_0+\eta L_{\log})$ descendants, then counting all joint
cells and using the parent maximal weight bounds the total mass of
such parents by $O(\exp(-(\eta-2\varepsilon)L_{\log}))$.
Parents of weight below $\exp(-h_0-\eta L_{\log})$ have total mass
at most $\exp(-(\eta-\varepsilon)L_{\log})$.  Outside these two
exceptions, the conditional support and maximal weight bounds are
the asserted ones, with errors $O((\eta+\varepsilon)L_{\log})$.
Choose $\eta\downarrow0$ much more slowly than $\varepsilon$.
The discarded fractions can thereby be made negligible compared
with any finite list of subpower densities needed later.

The support bound also shows that bins of absolute weight below
$\exp(-h-\eta L_{\log})$ have negligible total mass.  For a finite
list of joint labels, one may remove these bins successively, or
iteratively remove bins whose remaining weight is below the chosen
threshold.  Charge each removed bin when it is first deleted.  The
sum of all charges is bounded by the sum of the corresponding
support counts times their thresholds, and is still negligible if
the slack dominates the existing errors.  Dividing joint weights by
parent weights gives the corresponding conditional lower-weight
statements.

After any further restriction of density $\rho=\exp(-o(L_{\log}))$,
the support bounds remain valid and maximal weights grow by at most
$\rho^{-1}$.  Thus the absolute entropies change by $o(L_{\log})$.
The same parent and base-fiber trims recover the conditional
statements.  This proves heredity.  Finally apply the construction on
larger finite grids, with their size increasing sufficiently slowly.
Off-grid tests are sandwiched by admissible finer and coarser grid
states.  The domain has interior shift vectors with every coordinate
positive, and others with every coordinate negative, so such
sandwiches exist with arbitrarily small offset errors.  The
Lipschitz estimates of Lemma~\ref{ent:finite-states} complete the
diagonal extension.
\end{proof}

\subsection{The calibrated profile and conditional axis prediction}

\begin{proposition}[Calibrated entropy profile]\label{ent:states}
After the regularization of Lemma~\ref{ent:regularization} and passage
to a subsequence, the functions
\[
\frac{\Ent(Z(z))-\Ent(Z(0))}{L_{\log}}
\]
converge locally uniformly on $\Omega$ to a continuous Lipschitz
function $F$.  Each scalar partial lies in $[0,1]$ where defined.
The function is coordinatewise nondecreasing, and its off-diagonal
distributional second partials are nonpositive in the interior.
The same assertions hold for disjoint groups of coordinates advanced
together with nonnegative speeds, on every valid comparison
rectangle.  At full alignment,
\begin{equation}\label{eq:source-11}
 F(C,B,D,A,Y,R)\geq k_C C+k_B B+q_*Y,
 \qquad k_C=1+\kappa,\quad k_B=1-\kappa,
\end{equation}
with equality on
\[
\gamma(Y)=(-aY,-bY,(1-a)Y,(1-b)Y,Y,\lambda Y).
\]
\end{proposition}

\begin{proof}
The finite-state comparisons give equicontinuity, after reference
subtraction, with errors tending to zero in the present log unit.
One first takes a subsequence on a countable dense grid, and then
uses the admissible coordinatewise joins and bounded comparable
shifts to obtain local uniform convergence on the full domain.
Monotonicity and Lipschitz bounds pass to the limit.  The entropy
rectangle inequalities give submodularity.  Equivalently,
$-\partial_{ij}F$ is a nonnegative distribution, and hence a
nonnegative Radon measure, for $i\ne j$ in the interior.
Pinning coordinates or summing disjoint nonnegative coordinate
increments preserves those inequalities.  Comparisons on faces
follow by translating the valid rectangle into the interior and
passing to the limit; convexity of the domain makes these
translations available.

To verify the normalization in \eqref{eq:source-11}, put
\[
B_0=-\log K+
 \bigl(2x_0+(1-\kappa)d_0+q_*y_0\bigr)\log M.
\]
At full alignment, a state atom is contained in a tested ellipse
cell, with its displayed time bin and axis.  The simultaneous cap
from Section~\ref{sec:geo} therefore gives
\[
\Ent(Z(z))\geq
B_0+L_{\log}(k_C C+k_B B+q_*Y)-o(L_{\log}).
\]
For a nonmesh time use a nearby earlier mesh time, with the rebasing
and width errors controlled by \eqref{eq:source-8}.
On the calibrated path, Corollary~\ref{geo:heavy-count} supplies the
opposite support estimate.  Given an assigned high box and its time
bin, the reference-axis approximation at path aspect accuracy leaves
only $\exp(o(L_{\log}))$ possibilities for the labels
\eqref{eq:source-9}.  The same is true of the slight width
mismatches and time rebasing.  Thus the support logarithm is at most
the displayed lower bound plus $o(L_{\log})$.  In particular,
$\Ent(Z(0))=B_0+o(L_{\log})$.  Subtract this equality and divide by
$L_{\log}$.  Equality first holds on the fine mesh and then
everywhere on the line by continuity.  All the cap and list
estimates survive the subpower restriction used to regularize.
\end{proof}

We also track the referenced conditional entropy of $(t,r)$ given
$\mathcal B$ on the same realizing sequence.  Its limiting profile
will be denoted by $\Psi(Y,R)$.  The time-only profile is $sY$, since
the upper weights and upper support counts in \eqref{eq:source-1}
give both entropy bounds, with errors negligible in the present
unit.

\begin{lemma}[Conditional prediction of the axis]\label{ent:predictability}
For $\lambda>0$,
$\Psi(Y,R)=sY$ on $R\leq\lambda Y$.
For $\lambda<0$, this equality holds for every finite $(Y,R)$.
Parameter-coordinate increments of $F$ dominate those of $\Psi$.
In particular, if $\lambda<0$, then $\partial_YF\geq s$.

In the latter case set $c=b-a=-\lambda>0$.  In the region
\begin{equation}\label{eq:source-12}
\begin{gathered}
C>aR/c,\qquad B>bR/c,\qquad Y>-R/c,\\
D>(a-1)R/c,\qquad A>(b-1)R/c
\end{gathered}
\end{equation}
inside $\Omega$, $F$ is locally independent of $R$.
\end{lemma}

\begin{proof}
Given $P$ and a time prefix at a local epoch, all particles in that
fiber have the same measured line state at that epoch.  Their
assigned high boxes therefore pass through that state.  The
pointwise exclusion in Corollary~\ref{geo:angular-exclusion} applies to
\emph{all} high prelaw boxes, irrespective of subsequent restrictions
or additional labels.  It implies that the assigned axes in the
fiber have only subpower many possibilities at path aspect
precision.  Corollary~\ref{geo:axis-prediction} transfers this conclusion
to the common reference axis.  The conditional entropy excess of
the axis over the time prefix is consequently $o(L_{\log})$ at that
precision.

If $\lambda>0$, the epoch of offset $Y$ predicts the reference axis
at offset depth $\lambda Y$, hence also at every $R\leq\lambda Y$.
If $\lambda<0$, for any fixed $(Y,R)$ choose an earlier fixed epoch
$v\leq Y$ with $\lambda v>R$.  Its time prefix is determined by the
one at $Y$, and its aspect precision is finer than the requested
axis cutoff.  The common reference epoch is still earlier, on the
deeper-aspect end of the stretched interval.  Chaining
\eqref{eq:source-8} from that reference to $v$ loses only the worst
aspect precision, namely the precision at $v$.  Thus this epoch
also predicts the common reference axis to the requested accuracy.
This proves the asserted zero excess entropy, including its
absolute $o(L_{\log})$ bound before reference subtraction.
Lemma~\ref{ent:finite-states} proves the increment comparison with
$F$, and hence $\partial_YF\geq s$ in the decreasing-aspect case.

We give the spatial prediction argument for
\eqref{eq:source-12} with all its buffers.  Fix a state in that open
region and choose a sufficiently small $\eta>0$.  Test the path epoch
\[
v_*=-\frac{R+\eta}{c}.
\]
Two particles with the same displayed state have a common truncated
slope $r$ and time $t_Y=T_{y_0+\delta Y}$.  Their differences satisfy
\begin{align*}
|\Delta U_1|&\lesssim M^{-x_0-\delta C},&
|\Delta U_2-r\Delta U_1|
 &\lesssim M^{-x_0-d_0-\delta B},\\
|\Delta(X_{t_Y})_1|&\lesssim M^{-x_0-y_0-\delta D},&
|\Delta(X_{t_Y})_2-r\Delta(X_{t_Y})_1|
 &\lesssim M^{-x_0-y_0-d_0-\delta A}.
\end{align*}
Since $Y>v_*$, they also have the same epoch-$v_*$ prefix.  Choose
the assigned axis $r_*$ of either particle at that epoch.  Its
distance from $r$ is at most
$M^{-d_0-\delta R+o(\delta)}$: compare first with that particle's
reference axis at depth $R+\eta$, and then with its common prefix
$r$ at depth $R$.

In the $r_*$ frame, the velocity errors have the two offset depths
\[
C,\qquad \min(B,C+R),
\]
with absolute background depths $x_0$ and $x_0+d_0$.  If
$t_*=T_{y_0+\delta v_*}$, then
$|t_*-t_Y|\lesssim M^{-y_0-\delta v_*}$.
Using $X_{t_*}=X_{t_Y}+(t_*-t_Y)U$, the position errors in that
frame have offset depths
\[
\min(D,C+v_*),\qquad
\min(A,D+R,B+v_*,C+R+v_*),
\]
with background depths $x_0+y_0$ and $x_0+y_0+d_0$.
The path targets at this epoch are
\[
-av_*,\quad -bv_*,\quad (1-a)v_*,\quad (1-b)v_*.
\]
Every displayed error depth is strictly finer than its target for
small $\eta$.  Indeed this follows from the five inequalities in
\eqref{eq:source-12}; the additional comparisons
$C+R>bR/c$ and $D+R>(b-1)R/c$ follow from $c=b-a$.

Thus peers in this state lie in bounded enlargements of one
another's assigned epoch boxes.  The all-high-box angular exclusion
forces their assigned axes to agree up to a subpower factor at
offset depth $R+\eta$.  Comparison with each particle's reference
axis gives the same conclusion for the finer reference-axis bins.
The state therefore determines this extra axis precision with
entropy cost $o(L_{\log})$.  After refining the parameter, its
spatial labels at the original cutoffs can be updated with bounded
ambiguity by \eqref{eq:source-10}.  Hence increasing $R$ a little
at fixed other offsets costs zero in the limiting profile.  The
strict margins hold throughout a neighborhood, proving local
independence.
\end{proof}

\paragraph{The stationary reference axis as a conditioning label.}
When $\lambda=0$, let $R_0$ be the reference slope truncated at
absolute depth $d_0$; take a constant label if $d_0=0$.  Retain
$R_0$ and measure the spatial coordinates in its frame.  Normalize
the velocity widths by $M^{-x_0},M^{-(x_0+d_0)}$ and the position
widths by the same factors times $M^{-y_0}$.  Use common offset
depths $u$ for both velocity components, $z$ for both position
components, and $Y$ for time.  The domain is $z\leq u+Y$.

Write $W_\zeta$ for these spatial and time bins at
$\zeta=(u,z,Y)$, and $\widehat Z_\zeta=(R_0,W_\zeta)$ for the
recorded state.  We use the referenced conditional entropy
\begin{equation}\label{ent:stationary-axis-cancellation}
\frac{\Ent(W_\zeta\mid R_0)-\Ent(W_0\mid R_0)}{L_{\log}}
=\frac{\Ent(\widehat Z_\zeta)-\Ent(\widehat Z_0)}{L_{\log}}.
\end{equation}
This is an exact finite-stage identity.  The size of $\Ent(R_0)$
is unrestricted: it cancels, rather than being charged as a
subpower loss.  Include $R_0$ and the needed joint labels in the
regularization list.  Thus we work with conditional blocks of a
recorded label, without selecting a single axis bin at a supposed
subpower cost.  The resulting profile has the following properties.

\begin{proposition}[Stationary-aspect profile]\label{ent:stationary-profile}
There is a hereditary Lipschitz submodular profile $F(u,z,Y)$ on
$z\leq u+Y$, with scalar-group Lipschitz bounds $2,2,1$, satisfying
at alignment
\begin{equation}\label{eq:source-13}
\begin{gathered}
F(u,u+Y,Y)\geq2u+q_*Y,\\
F(-aY,(1-a)Y,Y)=-2aY+q_*Y,\qquad \partial_YF\geq s.
\end{gathered}
\end{equation}
At alignment, arbitrary ellipse tests in the normalized frame are
also available.  For their two width depths $l_1,l_2$, the width
term $2u$ is replaced by
\[
(1+\kappa)\min(l_1,l_2)+(1-\kappa)\max(l_1,l_2).
\]
Any additional axis label required for such a lower-bound test may
be included on the measured side.
\end{proposition}

\begin{proof}
\textbf{Conditional capacity and the heavy list.}
The finite-state and regularization arguments apply with the two
spatial components grouped and with $R_0$ recorded.  Conditional
entropy inherits the same comparison inequalities.  To check the
baseline, at alignment set
\[
b_\zeta=B_0+L_{\log}(2u+q_*Y),
\]
with $B_0$ as in the proof of Proposition~\ref{ent:states}.
The joint cap gives maximal atom weight
$\exp(-b_\zeta+o(L_{\log}))$ for $\widehat Z_\zeta$.
On a positive fiber $R_0=r$ of probability $p_r$, normalization
therefore gives
\[
\Ent(W_\zeta\mid R_0=r)
\ge b_\zeta+\log p_r-o(L_{\log}).
\]
Averaging over $r$ gives the lower bound
$b_\zeta-\Ent(R_0)-o(L_{\log})$.  On the calibrated path,
the joint high-box list bounds
$\Ent(\widehat Z_\zeta)\le b_\zeta+o(L_{\log})$.
Indeed, an assigned high box and time bin leave only
$\exp(o(L_{\log}))$ possibilities for the depth-$d_0$ reference-axis
bin, by Lemma~\ref{geo:calibration} with $\lambda=0$.
Consequently
\[
\Ent(W_\zeta\mid R_0)
=b_\zeta-\Ent(R_0)+o(L_{\log}).
\]
The upper estimate is an averaged conditional entropy statement;
no uniform support bound in every untrimmed axis fiber is needed.
The common $\Ent(R_0)$ term cancels in
\eqref{ent:stationary-axis-cancellation}, giving the alignment
lower bound and calibrated equality in \eqref{eq:source-13}.
Typical-parent regularization gives the conditional support and
weight bounds within the recorded $(R_0,\text{parent})$ fibers,
so these comparisons have the required hereditary realizations.

\paragraph{Conditional time increments.}
For brevity put $t_Y=T_{y_0+\delta Y}$.  The same single-epoch
prediction as in Lemma~\ref{ent:predictability} gives
$\Ent(R_0\mid\mathcal B,t_Y)=o(L_{\log})$ at every fixed
local cutoff.  The exact chain rule reads
\[
\Ent(t_Y\mid\mathcal B,R_0)
=\Ent(t_Y\mid\mathcal B)
 +\Ent(R_0\mid\mathcal B,t_Y)-\Ent(R_0\mid\mathcal B).
\]
Subtract this identity at two cutoffs.  The last term cancels,
and the middle terms are negligible.  Thus the time increment
retains slope $s$, and the finite-state conditional comparison
gives $\partial_YF\ge s$.  No bound on the unconditional entropy
of $R_0$, or on $\Ent(R_0\mid\mathcal B)$, is asserted.

\paragraph{Elliptic tests in the normalized frame.}
For an ellipse in the normalized frame, apply
the submultiplicativity of $W$ under the fixed background linear
normalization, as in Section~\ref{sec:geo}.  Its weight in offset
units is the negative exponential of the stated width term.  Each
cell, with its axis label included if necessary, is contained in
one of these ellipse tests up to bounded covering, proving the
same entropy lower bound.
\end{proof}

\subsection{Realizing tangents and conditional blocks}

A tangent at a point $z_0$ of a profile is a locally uniform limit,
on the corresponding rescaled domain, of
\[
F_{z_0,h}(z)=\frac{F(z_0+hz)-F(z_0)}{h},\qquad h\downarrow0.
\]
We may take several profiles to their tangents simultaneously.
Centers may lie on an alignment face or on an affine sheet in that
face.  Constraints strict at the center disappear locally after
rescaling; binding linear constraints retain their exact tangent
form.

\begin{lemma}[Realization of tangents]\label{ent:tangent-realization}
Every tangent used below, and every finite succession of such
tangents at finitely many sites, has a realizing sequence of the
same sheared states, with hereditary support and weight estimates
in its final log unit.  Required common physical cutoffs at
different sites can be kept identical.  Exact alignment identities
are preserved, before bounded dyadic rounding, whenever the center
and the tested cutoffs lie on the corresponding alignment.
\end{lemma}

\begin{proof}
First fix a radius $h$ and a finite test grid for the desired
tangent.  Approximate the original profile on the corresponding
grid of absolute cutoffs to an error much smaller than $h$.
Choose the original stage sufficiently far out that
$hL_{\log}\to\infty$, and that every previous logarithmic error
and every probability-class loss is $o(hL_{\log})$.
Lemma~\ref{ent:regularization} can be imposed on this finite list
and on its joint and conditional labels.  Now pass to a diagonal
sequence as $h\downarrow0$ and the grid grows.  This realizes the
tangent in its new log unit.  Repeating the construction handles
iterated tangents.  A finite collection of sites is treated by
putting all its cutoffs in the same list.  Coordinates prescribed
to be physically common are represented by the same absolute
cutoff, rather than by separately rounded approximations.

Inside a common parent bin, subtract the known bin centers and use
its spatial and parameter widths as units.  Formula
\eqref{ent:shear-update} still applies.  For instance a time
increment changes position by its normalized scalar increment
times normalized velocity, with known center terms subtracted
exactly.  At an aligned triple the position unit equals the
velocity unit times the time unit.  The defining affine equalities
of these depths are preserved under a centered rescaling through
alignment.  Dyadic rounding changes the units only by bounded
factors.

For comparable parent and refinement states, the conditional
support and maximal-weight estimates now use their entropy
difference, by Lemma~\ref{ent:regularization}.  This applies as
well after a fine parameter bin is included in the parent tuple:
the old spatial labels can be rebased to that parameter bin with
bounded ambiguity.  Lower typical weights, when needed, are
obtained by the same trims.

To see the resulting product density, fix a typical parent and
let $\mathsf A,\mathsf B$ be the supported bin sets of two
prescribed refinement labels.  Let $\mathsf E\subset
\mathsf A\times\mathsf B$ consist of pairs with positive
conditional joint probability.  Use joint labels with the parent
where nesting has bounded ambiguity; this does not change the
exponents.  Write $L$ for the final logarithmic unit.  If the
two marginal entropy exponents are $\alpha,\beta$ and the joint
exponent is $\alpha+\beta$, regularization gives
\begin{gather*}
|\mathsf A|\leq e^{(\alpha+o(1))L},\qquad
|\mathsf B|\leq e^{(\beta+o(1))L},\\
\max_{(a,b)\in\mathsf E}\Pr(a,b\mid\text{parent})
\leq e^{-(\alpha+\beta-o(1))L}.
\end{gather*}
The conditional law has mass one, so
$|\mathsf E|\geq e^{(\alpha+\beta-o(1))L}$ and hence
\[
\frac{|\mathsf E|}{|\mathsf A|\,|\mathsf B|}
\geq e^{-o(L)}.
\]
This is subpower density of supported bin pairs, with no
independence assertion.  All these statements hold at every
fixed finite power grid, and hence on successively finer grids
by the same diagonal choice.
\end{proof}

In projection arguments we will use the actual realized profiles.
A full-alignment test is based in a full profile, or in a tangent
centered through full alignment at each intervening step.  A slice
with other coordinates moved to create slack need preserve only
the aligned triples it tests.

\subsection{Peeling at an affine sheet}

\begin{lemma}[Peeling]\label{ent:peeling}
Let a Lipschitz, coordinatewise nondecreasing, submodular profile be
defined on a convex domain given locally by linear inequalities,
and let $\Sigma$ be an affine sheet in one of its faces.  Suppose
a group of coordinates is held fixed by a tangent direction to
$\Sigma$ that strictly increases all the coordinates still
attached outside that group.  Coordinates separated at earlier
steps may vary arbitrarily in this direction.  At almost every
point of $\Sigma$, a first tangent splits that group as an
independent additive summand.  A finite succession of such
separations can be made in the same first tangent.

Suppose that, after these separations, each group varies on
$\Sigma$ by one proportional family of linear forms, and that the
nonzero forms of different groups are pairwise nonproportional.
Then each group's restriction to its sheet diagonal is affine.
In particular, a separated singleton with nonconstant form has a
constant full partial derivative throughout the tangent.
\end{lemma}

\begin{proof}
We prove the assertion with an interior mollification first.  Write
$\mu_{ij}=-\partial_{ij}F\geq0$ for $i\ne j$.
Choose a compact sheet patch, a smooth cutoff on a slightly larger
patch, and an offset compact set lying strictly inside the tangent
domain.  Translate the sheet patch by $hu$, with $u$ in that
offset set and $h>0$ small.  If $i$ belongs to the group being
separated and $w$ is the chosen sheet-tangent direction, then
$w_i=0$ and
\[
-\partial_{iw}F=\sum_{j\ne i}w_j\mu_{ij}.
\]
The integral of the left side on the translated patch against
the cutoff is bounded uniformly in $h$ and $u$: integrate by
parts in the sheet direction $w$ and use the bounded first
derivative $\partial_iF$.  Terms from previously separated
groups have bounded absolute integrals by the estimates at their
earlier separation steps, using symmetry of mixed derivatives.
All the remaining terms outside the current group are
nonnegative, with strictly positive coefficients $w_j$.
Consequently each such $\mu_{ij}$ has a uniformly bounded
integral on the translated sheet patch.  Use nested, slightly
larger patches at successive steps.  This establishes all the
required cross-pair bounds before any tangent is taken.

Average these bounds also over compact interior offsets $u$.
For the rescaled function $F_{x,h}$ the corresponding mixed
derivative has density $h\mu_{ij}(x+hu)$ in the offset variable.
Thus its averaged mass on an offset compact set, integrated over
sheet centers $x$, is $O(h)$.  These bounds are summable for
dyadic $h$.  By positivity, summability, and exhaustion by
countably many compact subsets, at almost every sheet center all
the cross-pair measures tend to zero on compact subsets in every
limiting dyadic tangent.

For an unmollified profile, the off-diagonal negative mixed
derivatives are Radon measures.  Mollify at a radius much smaller
than $h$ times the strict interior margin of the offset compact
set.  Enlarge that compact set slightly when testing the measure.
The same averaged bounds pass to the measures and give the
conclusion just obtained.  This argument never takes a Hessian
trace on a binding face.

Vanishing cross derivatives give additive separation on the
convex open tangent domain.  One way to see this is to mollify
locally again: the gradient in a given group is then constant
along each fiber of all the other coordinates.  Connected
convex fibers show that it depends only on that group's
coordinates.  The resulting conservative field has a primitive
on their convex projection.  Patching these primitives gives a
sum of group functions, up to a constant; Lipschitz continuity
extends the identity to the closed domain.

Apply Rademacher's theorem in intrinsic Euclidean coordinates on
the relative interior of $\Sigma$
\cite[Theorem~3.1]{Heinonen2005Lipschitz}.  At almost every sheet center
the restriction of the original profile to $\Sigma$ is differentiable.
Its locally uniform linear blowup shows that the restriction of any
tangent there is affine on the whole tangent sheet.  This uses
intrinsic sheet measure, not ambient almost-everywhere differentiability.
After the preceding splitting it is a sum of Lipschitz functions
of the distinct group forms.  Such a sum can be affine only if
each summand is affine along its form.  For completeness, fix
one form and apply successive finite differences in directions
annihilating each of the other forms but not the fixed one.
Additional differences annihilate the affine right side.  The
result says that a sufficiently high distributional derivative
of the chosen one-variable function vanishes.  It is a
polynomial, and its global Lipschitz bound makes its degree at
most one.  A zero group form is constant on the sheet and
requires no assertion beyond that.  For a singleton with
nonzero form, its projected coordinate runs through the real
line on the sheet, so its full summand is affine and its full
partial is constant.
\end{proof}

\begin{remark}\label{ent:peeling-forms}
A useful sufficient condition for successive peeling is that the
distinct nonzero group forms are positive on a common vector.
They generate a pointed polyhedral cone.  An extreme ray has an
exposing functional zero on that ray and strictly positive on
the other generating rays.  This is the required sheet
direction for that group.  Remove it and repeat.  Zero forms
can first be separated as the group held fixed by the common
positive direction.  Proportional forms must be combined
before this procedure.

The tangent choices can be made measurably when their rates are
integrated over a sheet.  The dyadic blowups, jointly for any
countable list of profiles, form a precompact family in a
separable metrizable space of locally uniform convergence.
On a patch of fixed binding constraints, choose nested balls
from countable nets that are visited infinitely often.  This
selects a limiting tangent and a subsequence measurably.
Tangential derivatives of the original restriction agree for
all choices at its differentiability points.
\end{remark}

\subsection{Signed envelopes along a path}

The next lemma separates the cost of the coordinates that decrease
along a path from the remaining coordinates. It constructs a tangent
of the full profile and an auxiliary function of the decreasing
coordinates, obtained from tangents of slices with the remaining
cutoffs fixed at finer levels. The auxiliary function gives lower
bounds for ordered increments and an upper comparison that is exact
at the starting point. In the applications, projection tests continue
to use actual entropy profiles and their realized tangents with the
required alignments preserved. The auxiliary function is used only
through the increment comparisons and overlap identities proved below.

\begin{lemma}[Signed envelope]\label{ent:signed-envelope}
Let $F_0$ be a Lipschitz, coordinatewise nondecreasing, submodular
profile on a convex polyhedral cutoff domain $\mathcal D$.  All
increment comparisons below are required to have valid corners.
Let
\[
\gamma(y)=(N(y),Q(y)),\qquad y\in J,
\]
be a Lipschitz path in a boundary sheet of $\mathcal D$.  Suppose
each coordinate of $N$ is strictly decreasing, with its speed
bounded above and below in magnitude, and each coordinate of $Q$
is nondecreasing.  Constant coordinates may belong to $Q$.
Assume
\[
(N(\alpha),Q(\beta))\in\mathcal D\qquad(\alpha\leq\beta).
\]
Assume also the following exterior-slack property.  On compact
subintervals of $J$, one can choose a vector $e$ with strictly
positive $Q$-coordinates such that replacing $Q(\beta)$ by
$Q(\beta)+\varepsilon e$ preserves admissibility and opens every
constraint involving $Q$ by a positive constant times
$\varepsilon$.  At a path point, fix a strictly exterior value
of $Q$ supplied by this property.  Define the pure tangent
domain $\mathcal P$ to be the tangent domain of that fixed
slice at the path's $N$-coordinates.  The inequalities involving
$Q$ are strict, so $\mathcal P$ is given by the active
inequalities involving only $N$ and is independent of the chosen
exterior value.  Assume that $\mathcal P$ is either full space or
the three-dimensional half-space
$\{(C,B,R)\in\R^3:B\leq C+R\}$; in the latter case the pure
group consists exactly of $N=(C,B,R)$.

At almost every $y_0\in J$, put
\[
v=-N'(y_0),\qquad w=Q'(y_0),\qquad
n(t)=-tv,\qquad q(t)=tw.
\]
Thus $v_i>0$ and $w_j\geq0$.  There are a first tangent $F$ of
$F_0$ at $\gamma(y_0)$ and a Lipschitz, nondecreasing,
submodular function $G$ on the pure tangent domain $\mathcal P$,
both referenced to vanish at the origin, with the following
properties.
\begin{enumerate}[label=\textup{(\roman*)}]
\item
The function $F-G$ is nondecreasing under valid ordered increases
of $N$ at fixed $Q$.  Moreover,
\[
F(n(\alpha),q(\beta))-G(n(\alpha))
\]
is independent of $\alpha$ whenever $\alpha\leq\beta$.

\item
On any compact test region, if $H$ is sufficiently large that
$n(-H)\geq N$ and the comparison tuples are valid, then
\begin{equation}\label{eq:source-14}
F(N,Q)\leq
G(N)+F(n(-H),Q)-G(n(-H))
=G(N)+F^+(Q).
\end{equation}
Equality holds at the origin.  There is also equality for pure
increments in the overlap region: if $Q\geq q(\beta)$ and
$N\leq N'$ satisfy $N_i>n_i(\beta)$ for every $i$, and all four
tuples obtained from $N,N'$ and $Q,q(\beta)$ are valid, then
\[
F(N',Q)-F(N,Q)=G(N')-G(N).
\]

\item
If $\mathcal P=\R^d$, then $G$ is linear on every strict ordering
region of the normalized depths $x_i=N_i/v_i$.  There is a number
$m$ such that $G(n(t))=-mt$, and in every ordering region
\[
\sum_i v_i\partial_iG=m.
\]
When $x_i$ is strictly largest, $\partial_iG$ is a constant
marginal rate; this rate is a global upper bound for that partial.

If instead $\mathcal P$ has the constraint $B\leq C+R$, the same
affine rate holds on the tied line.  Write the corresponding
speeds as $a,b,c>0$, so $b=a+c$.  The $C$-partial is constant on
\[
\frac Ca>\max\left\{\frac Bb,\frac Rc\right\},
\]
and the $R$-partial is constant on
\[
\frac Rc>\max\left\{\frac Ca,\frac Bb\right\}.
\]
Each constant is a global upper bound for its partial,
interpreted distributionally in the interior.
\end{enumerate}
The tangent and all genericity requirements can be imposed
simultaneously for countably many ancillary profiles and exterior
slices.
\end{lemma}

\begin{proof}
\textbf{Fixed exterior slices and their rates.}
For each fixed later parameter $\beta$, the function
$y\mapsto F_0(N(y),Q(\beta))$ is Lipschitz on $y<\beta$.
Choose a countable dense family of parameters containing all
dyadic partition endpoints used below.  At common
differentiability points of the corresponding restrictions, define
\[
r_\beta(y)=-\frac{d}{dy}F_0(N(y),Q(\beta)),\qquad
m(y)=\lim_{\beta\downarrow y}r_\beta(y),
\]
where the limit is taken through that countable family.
Submodularity and monotonicity of $Q$ show that $r_\beta(y)$
increases as $\beta$ decreases to $y$.  All these rates are
bounded, so $m$ is bounded and measurable.

\paragraph{Finite increments with a moving tied parameter.}
The Lipschitz integral formula
\cite[Theorem~3.3]{Heinonen2005Lipschitz} applies to each fixed
restriction.  Its bounded rates are locally integrable.  Choose
$y_0$ to be a Lebesgue point of $m$ and of every applicable
fixed-$\beta$ rate, using \cite[Theorem~1.6.19]{Tao2011Measure}
after a cutoff on a larger compact interval.  There are only
countably many fixed restrictions.  A shrinking interval with fixed
rescaled endpoints has average oscillation bounded by a fixed
multiple of a centered interval's average, so the same Lebesgue-point
conclusion applies below.  Fix
$\alpha_1<\alpha_2<\beta_0$, and put
\begin{align*}
I_h={}&F_0\bigl(N(y_0+h\alpha_1),Q(y_0+h\beta_0)\bigr)\\
&-F_0\bigl(N(y_0+h\alpha_2),Q(y_0+h\beta_0)\bigr).
\end{align*}
For each fixed $\beta_*>y_0$ in the countable family, and all
sufficiently small $h>0$, ordered increment comparisons give
\[
\int_{y_0+h\alpha_1}^{y_0+h\alpha_2}r_{\beta_*}(s)\,ds
\leq I_h\leq
\int_{y_0+h\alpha_1}^{y_0+h\alpha_2}m(s)\,ds.
\]
These are finite-increment comparisons on valid rectangles,
or their limits along the ordered path.  If the moving parameter
$y_0+h\beta_0$ is outside the countable family, approximate it
by parameters in that family still later than the entire
increment interval.  Lipschitz continuity of $F_0$ and $Q$
passes the bounds to the limit; no continuity of derivatives
in the conditioning value is required.  Divide by $h$, apply
Lebesgue differentiation, and then let $\beta_*\downarrow y_0$ to obtain
\[
\lim_{h\downarrow0}\frac{I_h}{h}
=m(y_0)(\alpha_2-\alpha_1).
\]
Differentiability of the path at $y_0$ and the Lipschitz bound
permit its replacement by the linearizations $n,q$.  Hence
every tangent at this point has pure tied rate $m(y_0)$ when
the pure cutoffs are strictly finer than the fixed $Q$-path
time.  Continuity includes the endpoint of the ordered region.

\paragraph{Approximating the trace by exterior slices.}
We next obtain the same rate from strictly exterior slices.
Work on a compact path interval with a fixed exterior vector
$e$.  For a dyadic partition of mesh $\delta_k$, set, on each
interval $I=[a_I,b_I]$,
\[
Q^I=Q(b_I)+\delta_k e,
\qquad
m_k(y)=-\frac{d}{dy}F_0(N(y),Q^I),\quad y\in I.
\]
The entire countable family of these constants is fixed before
choosing generic centers.  Ordered comparison gives
$0\leq m_k\leq m$.  We claim that $m_k\to m$ in $L^1$ on
compact subintervals.

Fix a coarser dyadic interval $[a,b]$.  For every fine
$I\subset[a,b]$, one has $Q^I\leq Q(b)+\delta_k e$, with all
comparison tuples valid.  Sum the corresponding pure increments
to obtain
\[
\int_a^b m_k(y)\,dy
\geq F_0(N(a),Q(b)+\delta_k e)
      -F_0(N(b),Q(b)+\delta_k e).
\]
First let the fine mesh tend to zero while keeping $[a,b]$
fixed.  Lipschitz continuity removes the exterior shift.
Then refine the coarse partition.  Its fixed-endpoint rates
converge almost everywhere to the defining monotone trace,
so dominated convergence makes the sum of its tied increments
tend to $\int m$.  The upper bound $m_k\leq m$ proves the
claim.  After a subsequence of partitions, $m_k(y)\to m(y)$
almost everywhere.  The fine limit preceding the coarse limit
is essential here; it also handles constant coordinates of $Q$
without continuity assumptions on derivatives in those
coordinates.

\paragraph{The pure limit and overlap equality.}
Choose $y_0$ outside all exceptional sets for the profile, the
countable exterior slices, and the restrictions used below.
Blow up $F_0$ at $(N(y_0),Q(y_0))$ and, simultaneously, each
applicable slice $N\mapsto F_0(N,Q^I)$ at $N(y_0)$.
Equicontinuity and a diagonal subsequence give a first tangent
$F$ and slice tangents $G_k$.  For each fixed exterior slice,
strict mixed slack makes its rescaled domain contain every
compact subset of $\mathcal P$ eventually.  Differentiability
along the tied path gives
\[
G_k(n(t))=-m_k(y_0)t.
\]
For fixed $k$, its exterior $Q^I$ is strictly finer than the
$Q$-coordinates of any fixed compact blowup test once the
blowup radius is small.  Submodularity therefore gives
\[
\partial_iF\geq\partial_iG_k,\qquad i\in N,
\]
in distributions on the interior.  Valid face comparisons
follow by interior translation and continuity.  Take a locally
uniform subsequential limit of the referenced $G_k$, denoted by
$G$.  Then
\[
\partial_iF\geq\partial_iG,
\qquad G(n(t))=-m(y_0)t.
\]
Thus $F-G$ is nondecreasing under valid ordered pure increments.
The finite-increment trace calculation makes its tied pure
increments vanish for $\alpha<\beta$, and continuity includes
$\alpha=\beta$.  This proves (i) and
\eqref{eq:source-14}, including equality at the origin.

For overlap equality, first fix $q(\beta)$.  A pure tuple
strictly finer than $n(\beta)$ is bounded above and below by
two tied tuples with path times strictly below $\beta$.
They are valid.  Monotonicity of $F-G$ and its equal values at
the tied tuples imply that $F(N,q(\beta))-G(N)$ is constant
throughout the comparison region.  The pure increment identity
therefore holds at $q(\beta)$.  Increasing $q(\beta)$ to $Q$
can only decrease a pure increment of $F$, whereas (i) bounds
it below by the corresponding increment of $G$.  These bounds
agree, proving (ii).

\paragraph{Order sectors of the pure profile.}
It remains to identify the pure slice tangents.  Suppose first
that $\mathcal P$ is full space.  Fix an exterior slice $Q^I$
and a nonempty proper coordinate subset $S$.  For each later
parameter $\beta$ in the countable dense family, define on the
local valid restrictions
\begin{align*}
r^S_\beta(y)&=-\frac{d}{dy}
 F_0(N_S(y),N_{S^c}(\beta),Q^I),\qquad y<\beta,\\
r_S(y)&=\lim_{\beta\downarrow y}r^S_\beta(y).
\end{align*}
As $\beta$ decreases, the complementary cutoffs become finer,
so submodularity makes these rates decrease.  Require
differentiability of all these fixed restrictions and take
$y_0$ to be a common Lebesgue point of their rates and traces,
for every exterior slice and subset.

For $\alpha_1<\alpha_2<\beta_0$, put
\begin{align*}
I_h^S={}&F_0(N_S(y_0+h\alpha_1),N_{S^c}(y_0+h\beta_0),Q^I)\\
&-F_0(N_S(y_0+h\alpha_2),N_{S^c}(y_0+h\beta_0),Q^I).
\end{align*}
For fixed $\beta_*>y_0$ in the dense family, sufficiently close
to $y_0$ that the fixed-complement restriction is valid near
$y_0$, and for sufficiently small $h$, the reversed
ordered-increment bounds are
\[
\int_{y_0+h\alpha_1}^{y_0+h\alpha_2}r_S(s)\,ds
\leq I_h^S\leq
\int_{y_0+h\alpha_1}^{y_0+h\alpha_2}r^S_{\beta_*}(s)\,ds.
\]
As before, approximate a moving complementary parameter by
dense parameters later than the increment interval, and use
Lipschitz continuity to pass these finite-increment bounds to
the limit.  Divide by $h$, let $h\downarrow0$, and then let
$\beta_*\downarrow y_0$.  Lebesgue differentiation gives
$I_h^S/h\to r_S(y_0)(\alpha_2-\alpha_1)$.  Thus the slice
tangent has constant tied-$S$ rate $r_S(y_0)$ whenever $S$
is strictly finer than its complement.  The exterior value
$Q^I$ stays fixed throughout this blowup.

The complement need not be tied: bound its normalized depths
above and below by tied, strictly coarser depths.  Full pure
space and the fixed exterior slack make all these corners
valid locally.  Ordered increments sandwich the tied-$S$
rate between the same two constants.  For the slice tangent
$G_k$, abbreviate its rate $r_S(y_0)$ by $r_S$.

Write $x_i=N_i/v_i$ and order them as
$x_{(1)}\leq\cdots\leq x_{(d)}$.  Start with all normalized
depths equal to $x_{(1)}$, increase the remaining coordinates
together, and freeze each as its specified depth is reached.
The tied-subset calculation yields
\[
G_k(N)=r_{[d]}x_{(1)}+
\sum_{j=2}^{d}r_{S_j}\bigl(x_{(j)}-x_{(j-1)}\bigr),
\qquad S_j=\{i:x_i\geq x_{(j)}\},
\]
where $r_{[d]}=m_k(y_0)$.  Repeated depths cause no ambiguity,
since their differences vanish.  This proves linearity on
strict ordering regions.  Translating all normalized depths
by the same amount proves the identity for the weighted sum
of partials.  If $x_i$ is strictly largest, its partial is the
singleton rate $r_{\{i\}}/v_i$.  At any other point, make all
complementary coordinates sufficiently coarser to put $i$ in
this finest sector.  Submodularity shows that the original
partial is no larger than this singleton rate.  The finitely
many coefficients are bounded, so these assertions pass to
the locally uniform limit $G$.

For the pure constraint $B\leq C+R$, the function $B-C-R$ along
the path is nonpositive and vanishes at the center. Its derivative
there is zero, giving $b=a+c$. The full tied trace and
the singleton traces for $C,R$ still have valid comparisons.
For example, if
$C/a>\max(B/b,R/c)$, tied complements
$(B,R)=(bT,cT)$ with $T<C/a$ are valid, because
$bT<C+cT$.  Tied complements above and below the given
normalized complement sandwich its singleton increments.
Nonnegative connecting segments remain in the domain by
convexity and validity of their endpoints.  Thus the
$C$-partial is constant in the finest sector.  Taking a
sufficiently coarse tied complement at any other interior
point gives the global upper bound.  The argument for $R$
is identical.  These statements pass to $G$ in distributions.

All exceptional sets arise from differentiability and Lebesgue
differentiation for countably many fixed restrictions.
Intersect their full-measure sets and diagonalize the
equicontinuous blowups.  This proves the simultaneous assertion.
\end{proof}

The pure auxiliary profile need not inherit projection statements.
At a center on the boundary of a strict order sector, a
pure-gradient lower bound for the actual profile passes
distributionally to the part of its tangent domain consisting of
directions into that sector.  If this sector cone meets the tangent
interior and peeling makes the relevant singleton partial constant
throughout the tangent, its value on that open portion is the same
constant used everywhere else.  The sector bound therefore bounds
that constant.  The transfer uses
interior ordered-increment inequalities and their continuous
extension to valid boundary rectangles.

\Needspace{4\baselineskip}
\begin{remark}[Boundary increments and inherited flatness]
\label{ent:signed-envelope-boundary}
At a binding constraint, overlap equality applies in particular
to a strictly positive pure direction with valid comparison
tuples.  If individual pure-gradient lower bounds hold and a
subsequently tested tangent has constant partials in these
coordinates, equality of that weighted sum forces equality in
each lower bound: every weight is strictly positive and every
excess over its lower bound is nonnegative.  First pass each
bound to the interior part of its sector cone and identify the
constant partial there.  A two-sided coordinate perturbation
at the boundary origin is unnecessary.

Flatness of the exterior slices on a common strict comparison
region passes to their tangents and then to $G$.  In particular,
the flatness in \eqref{eq:source-12} is inherited wherever its
offset conditions hold for those slices.  Pure lower profiles
are used only through the properties just proved; projection
tests continue to use actual entropy profiles.

The exterior-slack vectors have explicit choices in the
applications.  On the calibrated path with $\lambda>0$ and
$a>0$, the negative coordinates are $C$ and whichever of
$B,D$ strictly decrease.  Large positive shifts in $R,Y$,
with smaller positive shifts in the other nonnegative
coordinates, open the mixed constraints.  For $\lambda<0$,
the negatives are $R,B$, and also $C$ if $a>0$.  Positive
shifts in $Y,D,A$, successively smaller in that order, open
the mixed constraints; add a positive shift in $C$ when it
is static.  In the scalar domain $Z\leq X+y$ with decreasing
$X$, shift $y$ more than $Z$.

The tied-subset argument also applies with signs reversed on
a full domain along a line whose speeds are all strictly
positive.  It gives the piecewise-affine order slopes of
unconditional or conditional two-variable profiles used below.
\end{remark}

The final lemma turns local comparisons into one finite path. Its
almost-everywhere hypothesis is along every admissible path, as its
statement specifies.

\subsection{From local gains to a finite path}

\begin{lemma}[Local gains and path control]\label{ent:control}
Fix a starting state and a duration $T>0$, and consider paths
from that state with velocities in a compact convex box.
The states may be restricted to a closed linear half-space;
assume there is an available velocity
following its boundary.  Let the potential be continuous and
Lipschitz on the compact region reachable by these paths.
If the potential depends on time, include time as a state
coordinate with its prescribed unit velocity.
Suppose that along \emph{every} admissible Lipschitz path, at
almost every interior time, arbitrarily short admissible
straight steps increase the potential at a rate exceeding
$g_0+\varepsilon$, where $\varepsilon>0$ is fixed.  Then a
reachable terminal state has potential at least $g_0T$ greater
than its value at the prescribed start.

The local hypothesis is satisfied at a point if a tangent
there has a strict admissible finite-step comparison from the
tangent origin with the stated margin.  A sufficiently small
admissible perturbation of its velocity or endpoint is harmless
within a strict margin.
\end{lemma}

\begin{proof}
Include time among the state coordinates, so the terminal
face is part of the compact reachable set.  Call a state good
if it admits a positive-length straight step before the
terminal time with gain rate greater than
$g_0+\varepsilon/2$, and call all other states bad.  Goodness
is relatively open.  Indeed a strict comparison persists
under small perturbations.  If a step ends on the terminal
face or on the boundary of the half-space, shorten it
slightly; the comparison remains strict.  A step starting
on the half-space boundary either enters the interior or
follows the boundary, and remains admissible when its start
is moved slightly into the half-space.  For an interior
start, shortening or a small convex perturbation toward a
boundary-following velocity keeps the perturbed step
admissible.  Thus the bad set is closed; it includes the
terminal face.

By hypothesis, the bad set occupies zero time on every
admissible path.  Its shrinking neighborhoods consequently
have uniformly vanishing occupation time over all such
paths.  To prove this, otherwise choose shrinking radii and
paths with a fixed positive occupation time.  The paths have
a common Lipschitz bound.  A uniformly convergent subsequence
has an admissible limit, because the velocity box is compact
and convex and the half-space is closed.  Uniform convergence
shows that the limiting path spends at least that positive
time in every fixed neighborhood of the closed bad set.
Continuity from above of Lebesgue measure then gives positive
time in the bad set itself, a contradiction.

Fix a small neighborhood of the bad set.  Outside it, compactness
and openness of goodness give finitely many comparison steps,
each valid on an open set of starting states and with its length
bounded below by a positive number.  At a state outside the
neighborhood, take one of these good steps.  Inside the
neighborhood, take an arbitrary viable straight step of a
fixed sufficiently small length, clipped at terminal time.
Choose that length so that such a step stays in a slightly
larger neighborhood, using the common speed bound.  This
constructs a path reaching the terminal time in finitely many
steps.  The total time of its nongood steps is at most the
occupation time of the larger neighborhood, hence is
uniformly small.  Its good steps gain at least
$g_0+\varepsilon/2$ per unit time, whereas the Lipschitz
bound gives a uniform lower rate on the remaining steps.
Taking the neighborhood small enough makes the total gain
at least $g_0$ times the full duration.

Finally, a strict endpoint comparison from the tangent origin
passes back along its defining blowup sequence: the endpoint error
is little-oh of the step length, while the gain margin is
a fixed positive multiple of that length.  Admissibility under
the linear domain constraint is retained in this passage.
Any additional perturbation must remain admissible and have
velocity in the given box; sufficiently small such perturbations
preserve the strict margin.  This gives the last assertion.
\end{proof}
